\documentclass[12pt, a4paper]{article}

%=============== PAQUETES ===============
\usepackage[utf8]{inputenc}
\usepackage[T1]{fontenc}
\usepackage[spanish,es-tabla]{babel}
\usepackage{geometry}
\usepackage{amsmath}
\usepackage{graphicx}
\usepackage{circuitikz}
\usetikzlibrary{babel}
\usepackage{siunitx}
\usepackage{hyperref}
\usepackage{enumitem}

%=============== CONFIGURACIÓN ===============
\geometry{a4paper, margin=2.5cm}

\hypersetup{
    colorlinks=true,
    linkcolor=blue,
    urlcolor=cyan,
    pdftitle={Examen 1: Fundamentos de Circuitos DC},
}

% Configuración de unidades SI
\sisetup{
    output-decimal-marker = {.},
    per-mode = symbol
}

%=============== DOCUMENTO ===============
\begin{document}

\title{\textbf{Examen 1} \\ \large Fundamentos de Circuitos de Corriente Continua (DC)}
\author{Plan de Estudios de Electrónica}
\date{\today}
\maketitle

\vspace{0.5cm}

\section*{Problemas a Resolver}

\begin{enumerate}[leftmargin=*, label=\textbf{\arabic*.}]

    % Problema 1 (Originalmente Ejercicio 2)
    \item \textbf{Asociación de Resistencias:} Determine la resistencia equivalente ($R_{eq}$) vista desde los terminales $A$ y $B$ para el circuito mostrado en la Figura \ref{fig:ej1}.
    \begin{figure}[h!]
        \centering
        \begin{circuitikz}[scale=1]
            \draw (0,3) node[anchor=east] {A} to[short, o-*] (1,3)
                  to[R, l=$R_1$ (\SI{5}{\ohm})] (4,3)
                  to[short, *-] (4,4)
                  to[R, l=$R_2$ (\SI{20}{\ohm})] (7,4)
                  to[short, -*] (7,3);
            \draw (4,3) to[short] (4,2)
                  to[R, l=$R_3$ (\SI{20}{\ohm})] (7,2)
                  to[short] (7,3);
            \draw (1,3) to[short] (1,1)
                  to[R, l=$R_4$ (\SI{15}{\ohm})] (7,1)
                  to[short] (7,3);
            \draw (7,3) to[short] (8,3)
                  to[R, l=$R_5$ (\SI{2.5}{\ohm})] (8,0)
                  to[short, -o] (0,0) node[anchor=east] {B};
        \end{circuitikz}
        \caption{Circuito para el problema de reducción de resistencias.}
        \label{fig:ej1}
    \end{figure}

    % Problema 2 (Originalmente Ejercicio 7)
    \item \textbf{Análisis Nodal:} Considere el circuito de la Figura \ref{fig:ej2}. Formule el sistema de ecuaciones aplicando la LKC y determine las tensiones de los nodos principales $A$ ($v_A$) y $B$ ($v_B$).
    \begin{figure}[h!]
        \centering
        \begin{circuitikz}[scale=1.2]
            \draw (0,0) -- (6,0);
            \node[ground] at (3,0) {};
            \draw (0,0) to[I, l=\SI{5}{\ampere}] (0,3)
                  to[short, -*] (2,3) node[above] {Nodo A ($v_A$)}
                  to[R, l=\SI{2}{\ohm}, *-*] (2,0);
            \draw (2,3) to[R, l=\SI{2}{\ohm}] (6,3) node[above] {Nodo B ($v_B$)}
                  to[R, l=\SI{4}{\ohm}, *-*] (6,0);
        \end{circuitikz}
        \caption{Circuito para el análisis nodal.}
        \label{fig:ej2}
    \end{figure}

    % Problema 3 (Originalmente Ejercicio 8)
    \item \textbf{Análisis de Mallas:} Para el circuito plano mostrado en la Figura \ref{fig:ej3}, defina las corrientes de malla en sentido horario y resuelva el sistema matricial o algebraico necesario para encontrar la corriente que atraviesa el resistor central de \SI{5}{\ohm}.
    \begin{figure}[h!]
        \centering
        \begin{circuitikz}[scale=1.2]
            \draw (0,0) to[V, l=\SI{12}{\volt}] (0,3)
                  to[R, l=\SI{4}{\ohm}] (3,3)
                  to[R, l=\SI{5}{\ohm}, *-*] (3,0)
                  to[short] (0,0);
            \draw (3,3) to[R, l=\SI{2}{\ohm}] (6,3)
                  to[V, l_=\SI{6}{\volt}] (6,0)
                  to[short] (3,0);
        \end{circuitikz}
        \caption{Circuito para el análisis por mallas.}
        \label{fig:ej3}
    \end{figure}

\end{enumerate}

\end{document}